\subsection{Complete evaluation}
Both stages completed all 855 dataset--seed--depth tasks, yielding 16,245 model score records on 57 datasets. No planned dataset or repetition is removed. The sample spans 32--12,958 observations and 2--43 predictors. All five outer partitions of \texttt{lymphography} require the documented singleton-class inner-CV fallback. Table~\ref{tab:performance} reports the main models; complete mixture-grid and tuning records are released separately. Unless stated otherwise, differences below are percentage points of accuracy and intervals are unadjusted dataset-bootstrap intervals.

\subsection{Uniformly increasing sight}
At depth three, the two-sighted tree improves mean accuracy over standard CART by 2.76 points [0.67, 5.21], but the median paired improvement is only 0.13 points, with 29 wins, three ties, and 25 losses. The Holm-adjusted signed-rank $p$-value is 0.611. Thus the positive mean interval does not imply a broadly consistent paired advantage or rejection by the separate location test. At depth six the mean difference is 0.38 points [$-1.89$, 2.97], with $p_H=1.000$. With unrestricted depth it becomes $-5.14$ points [$-8.70$, $-1.86$], with $p_H=0.012$. The mean realized depths are 15.1 for the unrestricted two-sighted tree and 10.1 for CART; the former reaches depth 58 in one fit. Greater realized depth coincides with poorer test performance here, without identifying its cause.

The pure two-sighted 20-tree forest gains 2.76 points [1.42, 4.27] over the custom greedy forest at depth three ($p_H=0.012$). Its depth-six difference is $-0.13$ points [$-2.35$, 2.05], whereas the unrestricted difference is $-6.21$ points [$-9.59$, $-3.22$] ($p_H<0.001$). Figure~\ref{fig:depth} shows the depth dependence. Uniformly increasing local sight therefore does not reliably improve prediction when final capacity is relaxed. Its mean forest fitting-cost ratio remains 27.4--31.7 relative to the custom control; the two-sighted single tree costs 126.1--298.2 times the optimized-library CART fit. These cross-implementation ratios characterize the present scorer, not all non-greedy algorithms.

\subsection{Sparse two-sighted replacements}
The strongest depth-matched mixed result is shallow. Replacing ten of 100 CART members increases mean accuracy from 0.7263 to 0.7369: a paired gain of 1.07 points [0.52, 1.70], with $p_H=0.0033$. The median gain is 0.20 points, and the dataset counts are 31 wins, 15 ties, and 11 losses. Balanced accuracy increases by 1.12 points. Mean measured constituent cost rises from 0.075 to 0.868 s, a ratio of 11.6; the median dataset-paired cost ratio is 5.5. This is an accuracy gain under matched shallow capacity, not under equal computation.

The direction survives both specified sensitivities, with a smaller effect. Giving equal weight to the 48 known dataset families gives 0.82 points [0.32, 1.42], with family-adjusted $p_H=0.018$. On the 38 tasks outside the named synthetic/constructed set, the mean effect is 0.60 points [0.14, 1.15]. The latter interval is descriptive and unadjusted; neither sensitivity makes the convenience sample representative of general workloads.

At depths six and unlimited, CART-mixture mean gains shrink to 0.48 points [0.13, 0.88] and 0.34 points [0.08, 0.64]. Both median effects are zero and both primary-family signed-rank $p_H$ values are 0.273. These positive mean intervals and non-rejected signed-rank tests concern different estimands; they are not evidence of equivalence or contradictory decisions about one common hypothesis. Mean cost ratios increase to 16.2 and 23.8. The custom-greedy mixture shows the same attenuation, but its shallow mean cost ratio is only 4.1 because that greedy implementation is slower than library CART. Comparator choice materially changes the apparent tradeoff.

\subsection{Training-selected conventional baselines}
The tuned random forest has mean accuracy 0.7896 and balanced accuracy 0.7541. Relative to it, the CART mixture's accuracy differences are $-5.27$ points [$-8.06$, $-2.88$] at depth three, $-1.83$ points [$-3.49$, $-0.51$] at depth six, and $-1.13$ points [$-2.12$, $-0.29$] without a depth limit. Corresponding tuned-family $p_H$ values are below 0.001, 0.021, and 0.036. These are benchmark-average disadvantages, not losses on every dataset. In 285 training-only forest selections, depth three is chosen 115 times, depth six 83 times, and unrestricted depth 87 times. The shallow option can be useful, but fixing it for every task is not justified by these selections.

Mean random-forest refit time is 0.101 s, compared with 0.868, 1.527, and 2.787 s for the CART mixtures. However, mean total conventional-forest selection-plus-refit time is 3.090 s. Figure~\ref{fig:tradeoff} and Table~\ref{tab:performance} exclude that selection cost from their fitting axis. They consequently cannot establish end-to-end budget dominance: the conventional baselines are tuned and the mixtures are not, and timing scopes and implementations differ. The evidence does not support practical superiority of the mixed learner over the training-selected conventional alternative.
